% ---------------------------------------------------------------------
%  10-limitations.tex -- limits of the method and how to evaluate it
% ---------------------------------------------------------------------
\section{Limitations and How to Evaluate HPB}
\label{sec:limitations}

HPB comes out of years of running real products, but this paper does not
prove that it works better than what you use today. It offers a method, a
notation and tools, and explains why I believe they help. This section says
where HPB fits, where it does not, what can go wrong, and how to find out
whether it pays off in your own organization.

% =====================================================================
\subsection{Where HPB fits, and where it does not}
\label{sec:limitations-fit}

HPB earns its keep when three things are true at once: the product has
\emph{identifiable versions} that customers or devices actually run;
\emph{several pieces of work} are in flight on the same product; and people
need to know \emph{which version} contains a given change. That is the daily
reality of products that combine hardware, firmware and software, of
software shipped in discrete versions with long support, and of regulated
contexts, where HPB supports formal configuration management
\citep{ieee828,iso10007} without replacing it.

HPB adds little in four situations. With \emph{continuous deployment
without versions} \citep{humble2010cd}, every merge is its own release and
baselines lose their meaning; a reduced timeline with deployments as
baselines is possible, but I do not claim it beats a deployment log. When
a product has \emph{one change at a time}, the diagram is a straight line
and a changelog is enough. \emph{Before the first release} there is
nothing to fork from; early exploration can live as lab branches of a
provisional line. And \emph{portfolio decisions}, how to spread people and
money across many products, need aggregated views the method does not
define yet.

% =====================================================================
\subsection{Known limitations}
\label{sec:limitations-known}

\paragraph{No empirical evidence yet.}
I report no data showing that HPB improves planning accuracy, traceability,
quality or decision speed, and none should be read into this paper. The
benefits I describe are hypotheses that follow from the structure of the
method; \cref{sec:limitations-evaluate} explains how to test them.

\paragraph{Designer bias.}
I designed HPB from my own practice and built the prototype in that same
context. The method may reflect the needs and vocabulary of that context
more than I can see, and the comparison in \cref{sec:gaps} rates HPB by
design, not by observed use.

\paragraph{Scale.}
The notation is designed for one product line with tens of open branches.
Collapsed history and lane packing (\cref{fig:scale}) stretch that range,
but screen width and human attention are still limits. A line with
hundreds of open branches needs filters, and a filtered diagram no longer
shows everything at a glance.

\paragraph{Garbage in, garbage out.}
Every metric and every audit answer depends on three fields being right:
the fork baseline, the closing baseline and the type of each branch. The
fork baseline of issues is the weak spot: when nobody records the version
in which a defect was observed, the default quietly blames recent
releases. HPB reduces this risk by making someone accountable for these
fields at intake and by flagging inferred values; it does not remove it.

\paragraph{Culture.}
HPB makes decisions visible, including the decision not to do something
and the decision to ship late. Where that feels threatening, people comply
on the surface: they never discard, or they re-target silently. No method
fixes that by itself; HPB can only make the symptoms show in the diagram.

\paragraph{Deliberate simplifications.}
A branch has one fork baseline and closes into one baseline
(\cref{rule:origin,rule:closure}), so a fix that must reach several
maintained versions is one branch per version, linked as propagations,
much like cherry-picking \citep{chacon2014progit}: change logs stay exact,
but widely back-ported fixes multiply branches. A line has one parent, so
merging references or a family of variants on a shared platform is
recorded as annotations and sibling lines. Security, regulatory or
customer-specific work are attributes, not new lanes. And for hardware the
moment of ``release'' is not crisp; HPB fixes a default per baseline kind,
and an organization that uses another must say so, or its lead times will
not be comparable.

% =====================================================================
\subsection{How to evaluate HPB in your organization}
\label{sec:limitations-evaluate}

You do not need a research project to find out whether HPB helps you. You
need one product line, two release cycles and the discipline to compare
before and after honestly.
\begin{enumerate}[leftmargin=1.8em,itemsep=1pt,topsep=3pt]
  \item \textbf{Take a baseline.} Rebuild the last year of one line in
        derived mode (\cref{sec:implementation}) or in HPB Markup, flag
        every fork baseline you had to guess, and compute the metrics you
        can.
  \item \textbf{Ask ten questions first.} Write down ten questions your
        team gets asked (which version first contained this fix, what is
        still open for the next release), time the answers with your
        current tools and check them against your release records.
  \item \textbf{Run HPB for two release cycles} with the practices of
        \cref{sec:operating}, without changing your team process at the
        same time.
  \item \textbf{Ask an equivalent set of questions again} with the
        diagram, built from the same data, so HPB does not win just because
        its data is cleaner.
  \item \textbf{Look at outcomes and cost:} rule violations, gate slip with
        what was moved or discarded, age of open work, escaped defects, and
        the time spent keeping HPB up to date.
\end{enumerate}

\noindent Four hypotheses are worth testing, each with the result that
would count against HPB:
\begin{description}[leftmargin=0pt,itemsep=1pt,font=\normalfont\bfseries]
  \item[H1 Comprehension.] Questions about origin, parallel work and
        closure are answered more accurately, and not more slowly, with the
        diagram. \emph{Against:} no gain, or a gain paid with much longer
        answers.
  \item[H2 Discipline.] Violations of R1--R4 fall to a low, stable rate.
        \emph{Against:} silent closures and releases with unverified
        committed work that stay frequent.
  \item[H3 Decisions.] More slips and descopes are recorded as decisions,
        and fewer branches stay open forever. \emph{Against:} unchanged age
        profiles, descopes still found out after the release.
  \item[H4 Cost.] Keeping HPB up to date costs little compared with the
        meetings it uses. \emph{Against:} data entry the team judges out of
        proportion, or quiet abandonment after the trial.
\end{description}

\noindent With few releases a year, differences between baselines are
mostly noise, so look at trends; everything looks better in the first
month, so compare the first cycle with the last; and never use these
numbers to judge individuals. If your evaluation shows that HPB does
\emph{not} help in your context, I would especially like to hear about it.

% =====================================================================
\subsection{What comes next}
\label{sec:limitations-next}

Four directions follow from these limits: a \emph{portfolio view} that
aggregates many lines sharing platforms or teams without breaking
\cref{rule:origin,rule:closure}; \emph{automatic links} to version control
and build systems, so that release tags, build numbers and hardware
revision records create baselines and merges by themselves, with semantic
versioning \citep{semver} giving the size of a software baseline; an
\emph{interchange format}, for which HPB Markup is the first candidate; and
a \emph{versioned specification} of the notation, changed only on evidence
about readability, not on taste. The conclusion says how to contribute.
